% !TeX root = ../papers/nc-p3.tex
\section{剰余標数3における正則性直線の全幾何点}
\label{sec:rank-fifteen-residue-line}
\begin{theorem}\label{thm:rank-fifteen-residue-line}
剰余標数3の付値を持つ標数零の体上で、通常な階数15窓の正則性座標が
$\bar x=0$、$y$ は単元を満たすとする。$\bar y$ は任意でよい。
この窓は標準Koszul左右AS正則で、各始点の経路の直和は
四つの中心四次式の多項式環上で階数256の自由加群である
（基礎体は代数閉体に拡大して考える）。
\end{theorem}
\begin{proof}
剰余体の代数閉包で $i^2=-1$ とする。$\bar y=H^6$、
$s^2+s^{-2}=2H$ と選ぶと、$s(s^4+1)\ne0$ である。
$t=is$ により $K=0$、$H=(s^2+s^{-2})/2$ となる。
四次式の根とその相方を
\[
 n=\tfrac12\bigl(s+s^{-1}-i(s-s^{-1})\bigr),\qquad
 N=\tfrac12\bigl(s+s^{-1}+i(s-s^{-1})\bigr)
\]
と取れる。実際 $nN=H$、$n^2+N^2=2=-1=K^2-1$ である。
ゆえに $U=0$, $V=H^2$ とすれば $(x,y)=(0,H^6)$、
$v=H^2$ は剰余三次式 $v(x-4v)^2=y$ の根となる。
こうして直線上の全剰余点に対し、係数を一変数のLaurent多項式で表示できる。

\path{computations/verify_rank_fifteen_residue_line.py} はこの表示を
$\mathbb F_9[s]$ 上で厳密に計算する。$P=s^2+1$, $Q=s^2-1$ とすると
$4s$ 倍した八係数 $(a,b,L,M,A,B,C,D)$ は
\[
 (4s,-2s,P-iQ,P+iQ,iQ-P,iQ+P,iP-Q,iP+Q).
\]
巡回する各隣接像の全生成行をここから取る。
分数体上の経路次元は長さ0から5で $(1,4,10,20,35,56)$、
四次中心の文字別次元は $(3,2,0,0)$ である。
この消去で逆にした主元の分子の全因子を記録する。
文字零の三自由座標における $(1,3,5),(2,7,11)$ と文字一の二基底を
四つの中心式に選び、各始点で分母を払う。
記録した主元が単元なら、これらは真の中心式を始点ごとに単元倍したものである。

次に各始点で、始点の四中心式を課した経路双対を長さ13まで計算する。
$\mathbb F_9[s]$ を基底 $1,i$ の $\mathbb F_3[s]$ 加群として扱い、
各段で終端の二次関係と、長さ4では中心式も課す。
得られる多項式行列の列像の基底を $C$ とする。$C$ の行加群の基底を並べた
正方行列の行列式は、単数を除いて元の行列の最大非零小行列式の最大公約元である。
これは二つの基底変更がそれぞれの生成加群を保つこととCauchy--Binet公式から従う。
その根の外では核が特殊化と可換なので、核と末尾文字ごとの射影を次段へ渡せる。
全始点で、係数拡大の次数2で割った商次元列は
\[
 (1,4,10,20,31,40,44,40,31,20,10,4,1,0)
\]
である。主元因子と全段の行列式の因子を合わせ、$\mathbb F_9$ 上で分解する。

$s\mapsto us,\ u/s$（$u^4=1$）は $H^6$ を保つ。
従って同じ $\bar y$ の認証には一つの代表で足り、因子の全軌道が残った場合だけ
直接調べればよい。$s(s^4+1)=0$ を除くと完全に残る既約因子の軌道は
次数1が一つ、次数2が五つである。各代表の一根で、全四始点の中心次元と
四中心商の所定の次元列を有限体で再検算した。他の共役根は
$\mathbb F_9$ 上のFrobeniusで移る。超越的な剰余値には主元や行列式の零点はない。
最後に定理\ref{thm:rank-fifteen-quotient-residues}のモニックな逐次持上げを
選んだ剰余代表に適用する。同じ標数零の $(x,y)$ を持つ通常窓への移送と、
中心式の持上げ、全次数の自由性・Koszul性の判定により結論を得る。
\end{proof}
